\section{\SPMPC{} Formulation}

\subsection{Augmented Prediction Model}

The mobile-base and path-progress state is
\begin{equation}
  \xr =
  \begin{bmatrix}
    p_x & p_y & \theta & v & s & \omega
  \end{bmatrix}^{\top},
\end{equation}
where $(p_x,p_y)$ is the base position, $\theta$ is the heading, $v$ is the linear velocity, $\omega$ is the angular velocity, and $s$ is the path-progress variable.
The control input is
\begin{equation}
  \vu =
  \begin{bmatrix}
    a & \alpha & v_s
  \end{bmatrix}^{\top},
\end{equation}
where $a$ is linear acceleration, $\alpha=\dot{\omega}$ is angular acceleration, and $v_s=\dot{s}$ is virtual path-progress speed.
The base and progress dynamics are
\begin{equation}
\begin{aligned}
  \dot{p}_x &= v\cos\theta, &
  \dot{p}_y &= v\sin\theta, &
  \dot{\theta} &= \omega,\\
  \dot{v} &= a, &
  \dot{s} &= v_s, &
  \dot{\omega} &= \alpha .
\end{aligned}
\label{eq:method_base_progress}
\end{equation}

The liquid response is represented by a control-oriented second-order modal model.
For each horizontal direction $i\in\{x,y\}$,
\begin{equation}
  \ddot{\eta}_i
  +2\zeta_i\omega_i\dot{\eta}_i
  +\omega_i^2\eta_i
  =-\kappa_i a_i,
  \label{eq:method_slosh_dynamics}
\end{equation}
where $\eta_i$ and $\dot{\eta}_i$ are the modal displacement and velocity, and $\omega_i$, $\zeta_i$, and $\kappa_i$ are the equivalent natural frequency, damping ratio, and input gain.
For planar mobile-base motion, the excitation channel is approximated by
\begin{equation}
  a_x=a,\qquad a_y\approx v\omega .
  \label{eq:method_excitation}
\end{equation}
This approximation is used for online prediction rather than as a direct liquid-surface observer.
It follows the same control-oriented modeling principle used in low-order slosh estimation and anti-slosh planning studies \cite{A01_Leva2021,A02_Guagliumi2021}.

The liquid state is
\begin{equation}
  \xl =
  \begin{bmatrix}
    \eta_x & \dot{\eta}_x & \eta_y & \dot{\eta}_y
  \end{bmatrix}^{\top},
\end{equation}
and the augmented state is
\begin{equation}
\begin{aligned}
  \xaug =
  [&p_x,\ p_y,\ \theta,\ v,\ s,\ \omega,\\
   &\eta_x,\ \dot{\eta}_x,\ \eta_y,\ \dot{\eta}_y]^\top .
\end{aligned}
\label{eq:method_augmented_state}
\end{equation}
After discretization, the horizon model is written compactly as
\begin{equation}
  \xaug_{k+1}=f_d(\xaug_k,\vu_k).
  \label{eq:method_discrete_model}
\end{equation}
This state choice is the key distinction from smoothness-only local planning: the current residual liquid phase is propagated into future candidate controls.

\subsection{Model Proxy and Objective}

The internal model proxy for liquid response is denoted by $H_{\mathrm{model}}$.
A modal proxy is
\begin{equation}
  H_{\mathrm{model}}^{\mathrm{modal}}(t)
  =
  c_h\sqrt{\eta_x^2(t)+\eta_y^2(t)},
  \label{eq:method_h_modal}
\end{equation}
where $c_h$ maps modal displacement to an approximate height scale.
When the quasi-static turning contribution is enabled, the proxy becomes
\begin{equation}
  H_{\mathrm{model}}^{\mathrm{full}}(t)
  =
  H_{\mathrm{model}}^{\mathrm{modal}}(t)
  +\mathbb{I}_{\mathrm{para}}\frac{R^2\omega^2(t)}{4g},
  \label{eq:method_h_full}
\end{equation}
where $R$ is the container radius and $\mathbb{I}_{\mathrm{para}}$ indicates whether the parabolic term is included.
These quantities remain model proxies and are not treated as ground-truth free-surface measurements \cite{C01_Ferrari2026}.

Let the reference path be $\vr(s)=[r_x(s),r_y(s)]^\top$ with tangent angle $\psi(s)$.
The tangent and normal vectors are
\begin{equation}
\begin{aligned}
  \bm{t}(s) &= [\cos\psi(s),\ \sin\psi(s)]^\top,\\
  \bm{n}(s) &= [-\sin\psi(s),\ \cos\psi(s)]^\top .
\end{aligned}
\end{equation}
With $\bm{e}(s)=[p_x,\ p_y]^\top-\vr(s)$, the contouring and lag errors are
\begin{equation}
  \ec=\bm{n}(s)^\top\bm{e}(s),\qquad
  \el=\bm{t}(s)^\top\bm{e}(s).
  \label{eq:method_contour_lag}
\end{equation}

At each control cycle, \SPMPC{} solves
\begin{equation}
\begin{aligned}
  \min_{\xaug_{0:N},\vu_{0:N-1}} \quad
  & \sum_{k=0}^{N-1}\ell_k+\ell_N\\
  \mathrm{s.t.}\quad
  & \xaug_0=\hat{\xaug}(t),\\
  & \xaug_{k+1}=f_d(\xaug_k,\vu_k),\\
  & \xaug_k\in\mathcal{X},\quad \vu_k\in\mathcal{U}.
\end{aligned}
\label{eq:method_ocp}
\end{equation}
The stage cost is
\begin{equation}
\begin{aligned}
  \ell_k={}&
    q_c e_{\mathrm{c},k}^{2}
    +q_l e_{\mathrm{l},k}^{2}
    -q_p v_{s,k}
    +q_v(v_k-v_{\mathrm{ref},k})^{2}\\
  &+q_{v_s}(v_{s,k}-v_{\mathrm{ref},k})^{2}\\
  &+\|\vu_k\|_R^2
    +\|\vu_k-\vu_{k-1}\|_S^2\\
  &+\|\bm{x}_{\ell,k}\|_{Q_s}^{2}.
\end{aligned}
\label{eq:method_stage_cost}
\end{equation}
The terminal cost retains terminal path-error and slosh-state penalties.
The feasible sets encode base and control limits on $v$, $\omega$, $a$, $\alpha$, and $v_s$.
Only the first optimized base command is executed, and the problem is solved again at the next control cycle.

\subsection{Soft Reference Governor}

The Slosh-risk Reference Governor is a soft module placed before the \MPCC{} solve.
It evaluates candidate speed-reference scaling factors
\begin{equation}
  \mathcal{B}=\{\beta_0,\ldots,\beta_{M-1}\}\subseteq[\beta_{\min},1],
\end{equation}
where each candidate defines $v_\beta=\beta v_{\mathrm{ref}}^{\mathrm{nom}}$.
A short internal rollout gives a risk score
\begin{equation}
  R(\beta)=
  \max_{0\leq k\leq N_g}
  \frac{H_k(\beta)}{H_{\lim}}.
\end{equation}
The feasible candidate set is
\begin{equation}
  \mathcal{B}_{\mathrm{feas}}
  =
  \{\beta\in\mathcal{B}\mid R(\beta)\leq \rho\}.
\end{equation}
The raw scale is the largest feasible candidate when such a candidate exists, and $\beta_{\min}$ otherwise.
After rate limiting, the filtered value $\beta_f$ defines
\begin{equation}
  v_{\mathrm{ref}}^g
  =
  \min\{v_{\mathrm{ref}}^{\mathrm{nom}},
  \max(v_{\min},\beta_f v_{\mathrm{ref}}^{\mathrm{nom}})\}.
\end{equation}
The governed reference enters \eqref{eq:method_stage_cost}.
It does not directly replace the optimized base command, and its role is reference adaptation rather than hard safety enforcement.
